\begin{figure}[t]
\centering
\begin{tikzpicture}[node distance=7mm and 9mm, every node/.style={font=\small}]
\node[primarybox, text width=0.25\linewidth] (t0) at (0.0,0.0) {\textbf{Identity}\\[-1mm]{\scriptsize exact primal/reduced models}};
\node[secondarybox, text width=0.25\linewidth] (t1) at (4.6,0.0) {\textbf{Progress}\\[-1mm]{\scriptsize ascent and residual}};
\node[primarybox, text width=0.25\linewidth] (t2) at (9.2,0.0) {\textbf{Convergence}\\[-1mm]{\scriptsize KŁ finite length}};
\node[successbox, text width=0.25\linewidth] (t3) at (0.0,-2.2) {\textbf{Identification}\\[-1mm]{\scriptsize strict margins}};
\node[secondarybox, text width=0.25\linewidth] (t4) at (4.6,-2.2) {\textbf{Local map}\\[-1mm]{\scriptsize tangent Jacobian}};
\node[primarybox, text width=0.25\linewidth] (t5) at (9.2,-2.2) {\textbf{Screening}\\[-1mm]{\scriptsize certified radius}};
\draw[arrPrimary] (t0) -- (t1);
\draw[arrPrimary] (t1) -- (t2);
\draw[arrPrimary] (t2) -- (t3);
\draw[arrPrimary] (t3) -- (t4);
\draw[arrPrimary] (t3) -- (t5);
\end{tikzpicture}
\caption{Theorem dependency graph. Arrows indicate logical dependence, not that any target has already been proved.}
\label{fig:theory-graph}
\end{figure}
